\documentclass[11pt]{article}
\usepackage[margin=1in]{geometry}
\usepackage{amsmath,amssymb,amsthm,mathtools}
\usepackage{hyperref}

\newtheorem{theorem}{Theorem}[section]
\newtheorem{lemma}[theorem]{Lemma}
\newtheorem{proposition}[theorem]{Proposition}
\newtheorem{corollary}[theorem]{Corollary}
\theoremstyle{definition}
\newtheorem{convention}[theorem]{Convention}
\newtheorem{question}[theorem]{Question}
\theoremstyle{remark}
\newtheorem{remark}[theorem]{Remark}

\newcommand{\Z}{\mathbb{Z}}
\newcommand{\Q}{\mathbb{Q}}
\newcommand{\F}{\mathbb{F}}
\newcommand{\height}{\operatorname{ht}}
\newcommand{\Mult}[1]{#1\,\cdot}
\DeclareMathOperator{\Frac}{Frac}
\newcommand{\Cset}{\mathcal{C}}

\title{The nonvanishing lemma for ribbon shape weights is false,\\
and the true statement is a commuting-words theorem}
\author{Clio}
\date{16 September 2026}

\begin{document}
\maketitle

\begin{abstract}
Let $R_e^{W}$ add an $e$-ribbon to a partition with a weight depending on the whole shape
of the ribbon.  The classification of the commuting pairs $[R_e^{W},R_f^{\bar W}]=0$ in
\cite{Q149} rested on one unproved lemma: that no weight value may vanish.  \textbf{That
lemma is false.}  For every $e<f$ with $\gcd(e,f)\ge3$ there are commuting pairs in which
almost every weight value is zero --- at $(e,f)=(3,6)$ the operator $R_3^{W}$ that adds
\emph{only} the L-tromino $(2,1)$, and nothing else, commutes with an explicit $R_6^{\bar W}$
supported on two ribbon shapes out of $32$.

The repair is not a patch but a better theorem.  Write $\sigma\otimes\rho$ for the shape
weight of rank $e+f$ obtained by concatenating the occupancy words of $\sigma$ and $\rho$
across a free letter.  Then the \emph{entire} one-bead sector of the commutator says exactly
\[
   \bar\sigma \;=\; \sigma\otimes\rho \;=\; \rho\otimes\sigma ,
\]
that is, $\sigma$ and $\rho$ \textbf{commute} under $\otimes$; and the Lyndon--Schützenberger
phenomenon holds in this graded semigroup: two commuting weights are, up to scalars, powers
of a common weight $\tau$ of rank $d=\gcd(e,f)$.  The proof is a four-line Euclidean descent
and uses \emph{no} nonvanishing hypothesis whatever.  Everything then follows:
\begin{itemize}
\item the weight values vanish nowhere iff $\tau$ does;
\item for $d\le2$ the rank-$d$ two-bead relation forces $\tau$ nowhere zero, so
  Theorem~1.4 of \cite{Q149} is \textbf{unconditionally true} there --- in particular
  Corollary~1.5 (the coprime headline, and the decisive test at $(1,3)$) is now a theorem,
  with the $\gcd(e,f)=1$ case needing no two-bead input at all;
\item for $d\ge3$ it is \textbf{false} as stated, and the missing solutions are exactly the
  lifts of the rank-$d$ weights with zeros.
\end{itemize}
\end{abstract}

\tableofcontents

\section{What this paper does}

\subsection*{The claim it refutes}

\cite[Lemma 8.2]{Q149} (``\textsf{lem:nonvanish}'') asserts: if $[R_e^{W},R_f^{\bar W}]=0$
and $W\not\equiv0\not\equiv\bar W$ over an integral domain, then no value of $W$ or $\bar W$
is zero.  It was proved there for $e\le2$ and left as gap (G1) for $e\ge3$, with the remark
that it is ``the only gap in the classification''.  The brief for the present session asked
for a proof at $e\ge3$.

It is false.  The smallest counterexample is $(e,f)=(3,6)$ and it is small enough to state
in full (Theorem~\ref{thm:counter36}): let
\[
  W(\text{the ribbon }(2,1))=-1,\qquad W=0 \text{ on }(3)\text{ and }(1,1,1),
\]
\[
  \bar W(\,(2,3,1)\,)=+1,\qquad \bar W(\,(2,1,2,1)\,)=-1,\qquad
  \bar W=0 \text{ on the other } 30 \text{ shapes of size }6 .
\]
Then $[R_3^{W},R_6^{\bar W}]=0$.  Both operators are nonzero; indeed $R_3^{W}$ is the
operator that adds an L-tromino, and only an L-tromino.

\subsection*{The theorem that replaces it}

\begin{convention}[concatenation of shape weights]\label{conv:otimes}
A \emph{shape weight of rank $m\ge1$} over a commutative ring $F$ is a function
$\sigma:\{0,1\}^{m-1}\to F$.  For shape weights $\sigma_1,\sigma_2$ of ranks $m_1,m_2$
define $\sigma_1\otimes\sigma_2$, of rank $m_1+m_2$, by
\[
  (\sigma_1\otimes\sigma_2)(y)\;:=\;\sigma_1\bigl(y_{[1,m_1-1]}\bigr)\,
       \sigma_2\bigl(y_{[m_1+1,\,m_1+m_2-1]}\bigr),
  \qquad y\in\{0,1\}^{m_1+m_2-1} .
\]
The letter $y_{m_1}$ is ignored.  This is associative but not commutative; write
$\tau^{\otimes k}$ for the $k$-fold power.
\end{convention}

\begin{theorem}[structure; §\ref{sec:euclid}]\label{thm:structure0}
Let $\sigma_1,\sigma_2$ be nonzero shape weights of ranks $m_1,m_2$ over a field $F$ with
$\sigma_1\otimes\sigma_2=\sigma_2\otimes\sigma_1$, and put $d=\gcd(m_1,m_2)$.  Then there
are a nonzero shape weight $\tau$ of rank $d$ and $c_1,c_2\in F^{\times}$ with
$\sigma_i=c_i\,\tau^{\otimes(m_i/d)}$.
\end{theorem}

This is the Lyndon--Schützenberger theorem (``two words commute iff they are powers of a
common word'') for the semigroup of Convention~\ref{conv:otimes}, and it is proved by the
same Euclidean descent.  Its relevance is Proposition~\ref{prop:1B}: \emph{the one-bead
sector of the ribbon commutator is precisely the statement $\bar\sigma=\sigma\otimes\rho
=\rho\otimes\sigma$.}  So the $\gcd$ in \cite[Thm.~5.2]{Q149} is not an accident of the
answer; it is the $\gcd$ of the commuting-words theorem.

\begin{theorem}[classification; §\ref{sec:main}]\label{thm:main0}
Let $1\le e<f$, $d=\gcd(e,f)$, $K$ an integral domain with fraction field $F$, and let
$W,\bar W$ be shape weights of ranks $e,f$ over $K$ with $W\not\equiv0\not\equiv\bar W$.
Then $[R_e^{W},R_f^{\bar W}]=0$ if and only if there are a nonzero shape weight $\tau$ of
rank $d$ over $F$ and $c,\bar c\in F^{\times}$ with
\[
   \sigma \;=\; c\,\tau^{\otimes(e/d)},\qquad
   \bar\sigma \;=\; \bar c\,\tau^{\otimes(f/d)},
\]
and $\sigma$ satisfies the two-bead relation \eqref{eq:2Bstar} of rank $e$ (vacuous if
$e=1$).  Here $\sigma(u)=(-1)^{|u|}W(u)$ as in \cite[Conv.~2.1]{Q149}.
\end{theorem}

\begin{corollary}[the corrected nonvanishing statement; §\ref{sec:main}]\label{cor:main0}
In the situation of Theorem~\ref{thm:main0}, $W$ and $\bar W$ vanish nowhere if and only if
$\tau$ does.  Consequently \cite[Lemma 8.2]{Q149} is
\[
  \textbf{true for } \gcd(e,f)\le2, \qquad \textbf{false for } \gcd(e,f)\ge3 .
\]
\end{corollary}

So the headline of \cite{Q149} survives and is upgraded: for $\gcd(e,f)=1$ the argument
becomes unconditional, and --- unexpectedly --- needs no two-bead input at all
(Corollary~\ref{cor:coprime}).  What fails is the $\gcd\ge3$ half of the classification,
which was the half that was new.

\subsection*{Why the old route could not have worked}

The abandoned strategy was to prove that the zero set of $\sigma$ is empty by propagating
zeros through cylinders.  Three warnings had been recorded against it, and all three were
right.  The $2405.15848$ template (an integrable structure whose defining operator has a
kernel by construction) said that ``commutation $\Rightarrow$ nonvanishing'' is not a general
principle; the Hammersley--Clifford obstruction said that ``a multiplicative structure on
$\{0,1\}^n$ forces a product support'' is false without positivity; and Terwilliger's
classification of the same theorem shape keeps nonvanishing and rigidity as \emph{separate}
axioms rather than deriving one from the other.  The lesson is recorded in
§\ref{sec:postmortem}: the $e\le2$ proof closes not because of anything about ribbons but
because at $e\le2$ the rank-$d$ weight $\tau$ has at most two values and a square relation
between them.

\section{Setting and notation}\label{sec:setup}

We work throughout in the $\sigma$-coordinates of \cite[Conv.~2.1]{Q149}: for a shape weight
$W$ of rank $e$,
\[
   \sigma(u)\;=\;(-1)^{|u|}W(u),\qquad u\in\{0,1\}^{e-1},
\]
where $u=u_M(b,e)$ is the occupancy word of the bead move $b\mapsto b+e$,
$u_i=[\,b+i\in M\,]$, and $|u|=\height$.  Murnaghan--Nakayama is $\sigma\equiv\text{const}$.
We freely use, and do not reprove, the following from \cite{Q149}:

\begin{itemize}
\item \cite[Lemma 3.1]{Q149} (\textsf{lem:realise}): every window configuration may be
  prescribed freely by a partition.  All the words appearing below are therefore realised by
  actual pairs of Maya diagrams, and no free parameter below is free for a trivial reason.
\item \cite[Lemma 4.1]{Q149} (\textsf{lem:sectors}): the commutator has matrix elements only
  in the one-bead and two-bead sectors.
\item \cite[Lemma 4.2]{Q149} (\textsf{lem:onebead}): with
  $X=M|_{(b,b+g)}=A\,m_e\,V\,m_f\,D$, $g=e+f$, $C=A\,m_e\,V$, $B=V\,m_f\,D$, the one-bead
  matrix element is, up to a global sign,
  \begin{equation}\label{eq:1B}
     \sigma(D)\,\bar\sigma(C)-\sigma(A)\,\bar\sigma(B),
  \end{equation}
  with $A,D\in\{0,1\}^{e-1}$, $V\in\{0,1\}^{f-e-1}$ and $m_e,m_f\in\{0,1\}$ free and
  independent.
\item \cite[Lemma 4.3]{Q149} (\textsf{lem:twobead}): for $2\le e$ and $1\le\delta\le e-1$ the
  two-bead matrix element is, up to a global sign,
  \begin{equation}\label{eq:2B}
     \sigma(\pi\,0\,\pi')\,\bar\sigma(\pi'\,0\,\theta')
     -\sigma(\pi\,1\,\pi')\,\bar\sigma(\pi'\,1\,\theta'),
  \end{equation}
  with $\pi\in\{0,1\}^{\delta-1}$, $\pi'\in\{0,1\}^{e-\delta-1}$,
  $\theta'\in\{0,1\}^{\delta+f-e-1}$ free and independent; the marked letter sits at
  position $\delta$ of the $\sigma$-argument and at position $e-\delta$ of the
  $\bar\sigma$-argument.  For $e=1$ the sector is empty.
\end{itemize}

\begin{convention}\label{conv:2Bstar}
Let $\sigma$ be a shape weight of rank $m\ge2$.  Say that $\sigma$ satisfies
\textbf{(2B$^*_m$)} if for every $\delta$ with $1\le\delta\le m-1$ and all
$\pi,\varsigma\in\{0,1\}^{\delta-1}$ and $\pi'\in\{0,1\}^{m-\delta-1}$,
\begin{equation}\label{eq:2Bstar}
  \sigma(\pi\,0\,\pi')\,\sigma(\pi'\,0\,\varsigma)
  \;=\;\sigma(\pi\,1\,\pi')\,\sigma(\pi'\,1\,\varsigma),
\end{equation}
the marked letter sitting at position $\delta$ in the first factor and at position
$m-\delta$ in the second.  For $m=1$ the condition is vacuous.

Say that $\sigma$ satisfies the stronger \textbf{(2B$^{**}_m$)} if the same identity holds
with the second factor's prefix decoupled: for all $\alpha,\beta'\in\{0,1\}^{\delta-1}$ and
$\alpha',\beta\in\{0,1\}^{m-\delta-1}$,
\begin{equation}\label{eq:2Bstarstar}
  \sigma(\alpha\,0\,\alpha')\,\sigma(\beta\,0\,\beta')
  \;=\;\sigma(\alpha\,1\,\alpha')\,\sigma(\beta\,1\,\beta').
\end{equation}
Taking $\beta=\alpha'$, $\varsigma=\beta'$ shows (2B$^{**}_m$)$\Rightarrow$(2B$^*_m$).
\end{convention}

It is convenient to record \eqref{eq:2Bstar} in \emph{window form}.  Let
$X=(X_1,\dots,X_{\delta+m-1})\in\{0,1\}^{\delta+m-1}$ and put
$\Omega_c(X)=\sigma(X_{c+1},\dots,X_{c+m-1})$.  Then \eqref{eq:2Bstar} at $\delta$ says:

\begin{equation}\label{eq:2Bwindow}
  \Omega_0(X)\,\Omega_\delta(X) \ \text{ is unchanged when } (X_\delta,X_m)=(0,0)
  \text{ is replaced by } (1,1),
\end{equation}
all other letters of $X$ being arbitrary and held fixed.  (Indeed $\Omega_0=\sigma(\pi
X_\delta\pi')$ with $\pi=X_{[1,\delta-1]}$, $\pi'=X_{[\delta+1,m-1]}$, and
$\Omega_\delta=\sigma(\pi' X_m\varsigma)$ with $\varsigma=X_{[m+1,\delta+m-1]}$, the letter
$X_m$ sitting at relative position $m-\delta$.)

\section{Eliminating $\bar\sigma$: the one-bead sector is $\otimes$-commutation}
\label{sec:reduce}

\begin{proposition}\label{prop:1B}
Let $1\le e<f$ and let $\sigma,\bar\sigma$ be shape weights of ranks $e,f$ over an integral
domain $K$ with fraction field $F$, both $\not\equiv0$.  Then \eqref{eq:1B} vanishes for all
admissible arguments if and only if there is a shape weight $\rho$ of rank $f-e$ over $F$
with
\begin{equation}\label{eq:comm}
   \bar\sigma \;=\; \sigma\otimes\rho \;=\;\rho\otimes\sigma .
\end{equation}
In that case $\rho$ is unique and $\rho\not\equiv0$.
\end{proposition}

\begin{proof}
Write \eqref{eq:1B}$=0$ as
\begin{equation}\label{eq:1Beq}
  \sigma(D)\,\bar\sigma(A\,m_e\,V)\;=\;\sigma(A)\,\bar\sigma(V\,m_f\,D)
  \qquad (A,D\in\{0,1\}^{e-1},\ V\in\{0,1\}^{f-e-1},\ m_e,m_f\in\{0,1\}),
\end{equation}
legitimate because $C=A m_e V$ and $B=V m_f D$ and all five blocks are free
(\cite[Lemma 3.1]{Q149}).

$(\Leftarrow)$ If \eqref{eq:comm} holds then $\bar\sigma(A m_e V)=\sigma(A)\rho(V)$ by the
first equality and $\bar\sigma(V m_f D)=\rho(V)\sigma(D)$ by the second, so both sides of
\eqref{eq:1Beq} equal $\sigma(A)\sigma(D)\rho(V)$.

$(\Rightarrow)$ Fix once and for all a word $A_0\in\{0,1\}^{e-1}$ with
$s_0:=\sigma(A_0)\neq0$; it exists because $\sigma\not\equiv0$.  \emph{Every cancellation
below is by this one element $s_0$.}
Specialising \eqref{eq:1Beq} at $A=A_0$ gives
\begin{equation}\label{eq:dagger}
  \sigma(D)\,\bar\sigma(A_0\,m_e\,V)\;=\;s_0\,\bar\sigma(V\,m_f\,D) ,
\end{equation}
and at $D=A_0$ gives
\begin{equation}\label{eq:ddagger}
  s_0\,\bar\sigma(A\,m_e\,V)\;=\;\sigma(A)\,\bar\sigma(V\,m_f\,A_0).
\end{equation}
Putting $D=A_0$ in \eqref{eq:dagger} and cancelling $s_0$ (legitimate: $K$ is a domain and
$s_0\ne0$):
\[
   \bar\sigma(A_0\,m_e\,V)\;=\;\bar\sigma(V\,m_f\,A_0)\qquad\text{for all }m_e,m_f,V .
\]
The left-hand side does not involve $m_f$ and the right-hand side does not involve $m_e$;
since they are equal, both are independent of both letters.  Define
\[
  \rho(V)\;:=\;\bar\sigma(A_0\,0\,V)/s_0\ \in F,\qquad V\in\{0,1\}^{f-e-1},
\]
so that $\bar\sigma(A_0\,m_e\,V)=s_0\rho(V)=\bar\sigma(V\,m_f\,A_0)$ for either value of
each marked letter.  Substituting into \eqref{eq:dagger} and cancelling $s_0$ gives
$\bar\sigma(V\,m_f\,D)=\rho(V)\sigma(D)$, i.e.\ $\bar\sigma=\rho\otimes\sigma$; substituting
into \eqref{eq:ddagger} and cancelling $s_0$ gives $\bar\sigma(A\,m_e\,V)=\sigma(A)\rho(V)$,
i.e.\ $\bar\sigma=\sigma\otimes\rho$.

Uniqueness is forced by $\rho(V)=\bar\sigma(A_0\,0\,V)/s_0$.  Finally $\rho\not\equiv0$,
since $\rho\equiv0$ would give $\bar\sigma\equiv0$.
\end{proof}

\begin{remark}
Proposition~\ref{prop:1B} supersedes \cite[Lemma 8.1]{Q149} (\textsf{lem:flip}).  The four
flip relations (F0)--(F3) proved there are all immediate consequences of \eqref{eq:comm}:
(F0) is the fact that $\sigma\otimes\rho$ ignores the letter in position $e$ and
$\rho\otimes\sigma$ the one in position $f-e$; (F1)--(F3) are the two factorisations written
out.  More to the point, Proposition~\ref{prop:1B} assumes \emph{no} nonvanishing.
\end{remark}

\begin{proposition}\label{prop:2B}
Assume the situation of Proposition~\ref{prop:1B} and $e\ge2$.  Then \eqref{eq:2B} vanishes
for all admissible arguments if and only if $\sigma$ satisfies \textnormal{(2B$^*_e$)}.
\end{proposition}

\begin{proof}
Fix $\delta$ and split $\theta'=\varsigma\,c\,\psi$ with $|\varsigma|=\delta-1$,
$c\in\{0,1\}$, $|\psi|=f-e-1$; this is possible since
$|\theta'|=\delta+f-e-1=(\delta-1)+1+(f-e-1)$.  The word $\pi'\,b\,\theta'
=\pi'\,b\,\varsigma\,c\,\psi$ of length $f-1$ has first $e-1$ letters $\pi'\,b\,\varsigma$,
letter $c$ in position $e$, and letters $\psi$ in positions $e+1,\dots,f-1$.  Hence
$\bar\sigma(\pi'b\theta')=(\sigma\otimes\rho)(\pi'b\theta')
=\sigma(\pi'\,b\,\varsigma)\,\rho(\psi)$, and \eqref{eq:2B} equals
\[
   \bigl[\sigma(\pi0\pi')\sigma(\pi'0\varsigma)-\sigma(\pi1\pi')\sigma(\pi'1\varsigma)\bigr]
   \cdot\rho(\psi).
\]
If (2B$^*_e$) holds this vanishes for every $\psi$.  Conversely, choose $\psi_0$ with
$\rho(\psi_0)\ne0$ (possible as $\rho\not\equiv0$) and cancel it; note this is a second,
\emph{separate} domain appeal, and its witness is $\rho(\psi_0)$, not $s_0$.  Since
$\pi,\pi',\varsigma$ are free, \eqref{eq:2Bstar} follows.
\end{proof}

\section{The Euclidean descent}\label{sec:euclid}

\begin{proof}[Proof of Theorem~\ref{thm:structure0}]
Induction on $m_1+m_2$.

\emph{Base $m_1=m_2=:m$.}  Then $d=m$, and a word $y\in\{0,1\}^{2m-1}$ has
$y_{[1,m-1]}=:u$ and $y_{[m+1,2m-1]}=:v$ free and independent, so the hypothesis reads
$\sigma_1(u)\sigma_2(v)=\sigma_2(u)\sigma_1(v)$ for all $u,v$.  Choose $u_0$ with
$\sigma_1(u_0)\ne0$ and set $\lambda=\sigma_2(u_0)/\sigma_1(u_0)$; taking $u=u_0$ gives
$\sigma_2=\lambda\sigma_1$, and $\lambda\ne0$ as $\sigma_2\ne0$.  Take $\tau=\sigma_1$,
$c_1=1$, $c_2=\lambda$.

\emph{Step.}  The hypothesis is symmetric in the two weights, so assume $m_1<m_2$ and put
$m_3=m_2-m_1\ge1$.  Parametrise $y\in\{0,1\}^{m_1+m_2-1}$ as
\[
   y \;=\; A\,\mu\,B\,\nu\,C,\qquad
   |A|=|C|=m_1-1,\quad |B|=m_3-1,\quad \mu,\nu\in\{0,1\},
\]
which is a partition of the $m_1+m_2-1$ letters because
$2(m_1-1)+2+(m_3-1)=m_1+m_2-1$.  Then $y_{[1,m_1-1]}=A$,
$y_{[m_1+1,\,m_1+m_2-1]}=B\,\nu\,C$, $y_{[1,m_2-1]}=A\,\mu\,B$ and
$y_{[m_2+1,\,m_1+m_2-1]}=C$, so the hypothesis $\sigma_1\otimes\sigma_2
=\sigma_2\otimes\sigma_1$ reads
\begin{equation}\label{eq:star}
   \sigma_1(A)\,\sigma_2(B\,\nu\,C)\;=\;\sigma_2(A\,\mu\,B)\,\sigma_1(C)
   \qquad\text{for all }A,\mu,B,\nu,C .
\end{equation}
Fix $A_0$ with $s_0:=\sigma_1(A_0)\ne0$; again \emph{every} cancellation below is by this
one element.
\begin{enumerate}
\item[(i)] Put $A=C=A_0$ in \eqref{eq:star} and cancel $s_0$:
  $\sigma_2(B\,\nu\,A_0)=\sigma_2(A_0\,\mu\,B)$ for all $B,\mu,\nu$.  The left side does not
  involve $\mu$ and the right side does not involve $\nu$, so both are independent of both.
  Define $\sigma_3(B):=\sigma_2(A_0\,0\,B)/s_0\in F$, a shape weight of rank $m_3$; then
  $\sigma_2(A_0\,\mu\,B)=\sigma_2(B\,\nu\,A_0)=s_0\,\sigma_3(B)$ for either marked letter.
\item[(ii)] Put $C=A_0$ in \eqref{eq:star}: $\sigma_1(A)\,\sigma_2(B\,\nu\,A_0)
  =\sigma_2(A\,\mu\,B)\,s_0$.  By (i) the left side is $\sigma_1(A)s_0\sigma_3(B)$; cancel
  $s_0$ to get $\sigma_2(A\,\mu\,B)=\sigma_1(A)\sigma_3(B)$, i.e.\
  $\sigma_2=\sigma_1\otimes\sigma_3$.
\item[(iii)] Put $A=A_0$ in \eqref{eq:star}: $s_0\,\sigma_2(B\,\nu\,C)
  =\sigma_2(A_0\,\mu\,B)\,\sigma_1(C)=s_0\sigma_3(B)\sigma_1(C)$; cancel $s_0$ to get
  $\sigma_2=\sigma_3\otimes\sigma_1$.
\end{enumerate}
Hence $\sigma_1\otimes\sigma_3=\sigma_2=\sigma_3\otimes\sigma_1$, and $\sigma_3\ne0$
because $\sigma_2\ne0$.  Since $\gcd(m_1,m_3)=\gcd(m_1,m_2)=d$ and
$m_1+m_3<m_1+m_2$, the inductive hypothesis gives a nonzero rank-$d$ weight $\tau$ and
$c_1,c_3\in F^\times$ with $\sigma_1=c_1\tau^{\otimes(m_1/d)}$ and
$\sigma_3=c_3\tau^{\otimes(m_3/d)}$.  By (ii) and associativity of $\otimes$,
\[
   \sigma_2\;=\;\sigma_1\otimes\sigma_3\;=\;c_1c_3\,
   \tau^{\otimes(m_1/d)}\otimes\tau^{\otimes(m_3/d)}
   \;=\;c_1c_3\,\tau^{\otimes(m_2/d)} . \qedhere
\]
\end{proof}

\begin{remark}[what the descent is]
Steps (i)--(iii) are the standard proof that two commuting words in a free monoid are powers
of a common word, transported to the graded semigroup of Convention~\ref{conv:otimes}.  The
statement ``$\sigma_2=\sigma_1\otimes\sigma_3=\sigma_3\otimes\sigma_1$'' is the assertion
that $\sigma_1$ is both a prefix and a suffix of $\sigma_2$ and that the two overlapping
occurrences agree --- the defect theorem.  No positivity, no nonvanishing, and no
invertibility of any weight value is used: the single witness $s_0$ does all the work.
\end{remark}

\section{The two-bead relation descends to rank $d$}\label{sec:descend}

\begin{lemma}\label{lem:descent}
Let $d\mid m$, $k=m/d$, let $\tau$ be a nonzero shape weight of rank $d$ over a domain, and
put $\sigma=\tau^{\otimes k}$ (rank $m$).
\begin{enumerate}
\item[\textup{(a)}] If $\tau$ satisfies \textnormal{(2B$^{**}_d$)} then $\sigma$ satisfies
  \textnormal{(2B$^*_m$)}.  (For $k=1$ this is trivial, and \textnormal{(2B$^{*}_d$)}
  suffices.)
\item[\textup{(b)}] If $\sigma$ satisfies \textnormal{(2B$^*_m$)} then $\tau$ satisfies
  \textnormal{(2B$^*_d$)}.
\end{enumerate}
\end{lemma}

\begin{proof}
Work in the window form \eqref{eq:2Bwindow}.  For $X\in\{0,1\}^{\delta+m-1}$ put
$T_c:=\tau(X_{c+1},\dots,X_{c+d-1})$.  Unwinding Convention~\ref{conv:otimes},
\[
   \Omega_c(X)\;=\;\sigma(X_{c+1},\dots,X_{c+m-1})\;=\;\prod_{i=0}^{k-1}T_{c+id},
\]
the sites $c+d,c+2d,\dots,c+(k-1)d$ being the ignored letters.  So
\[
   \Omega_0(X)\,\Omega_\delta(X)\;=\;\prod_{i=0}^{k-1}T_{id}\;\cdot\;
       \prod_{i=0}^{k-1}T_{\delta+id}.
\]
The factor $T_c$ involves the site $s$ iff $c<s<c+d$.  Now:
\begin{itemize}
\item \emph{site $\delta$.}  In the first product $c=id$, so we need $id<\delta<(i+1)d$:
  this happens iff $d\nmid\delta$, and then for the single index $p:=\lfloor\delta/d\rfloor$.
  In the second product $c=\delta+id\ge\delta$, so never.
\item \emph{site $m=kd$.}  In the first product we would need $id<kd<(i+1)d$ with
  $i\le k-1$: impossible.  In the second we need $i<(m-\delta)/d<i+1$, which happens iff
  $d\nmid\delta$, and then for the single index $i=k-p-1$; the corresponding factor is
  $T_{\delta+(k-p-1)d}=T_{(k-1)d+\delta_0}$, where $\delta_0:=\delta-pd\in[1,d-1]$.
\end{itemize}

\emph{Case $d\mid\delta$.}  No factor involves $X_\delta$ or $X_m$, so
\eqref{eq:2Bwindow} holds identically; nothing to prove in either direction.

\emph{Case $d\nmid\delta$, $\delta=pd+\delta_0$ with $1\le\delta_0\le d-1$.}  Let
$\Pi$ denote the product of the remaining $2k-2$ factors; $\Pi$ does not involve $X_\delta$
or $X_m$.  Writing $T^{(b)}$ for the value with $(X_\delta,X_m)=(b,b)$,
\eqref{eq:2Bwindow} at $\delta$ reads
\begin{equation}\label{eq:reduced}
   \Pi\cdot T^{(0)}_{pd}\,T^{(0)}_{(k-1)d+\delta_0}
   \;=\;\Pi\cdot T^{(1)}_{pd}\,T^{(1)}_{(k-1)d+\delta_0}.
\end{equation}
Explicitly, $T^{(b)}_{pd}=\tau(\alpha\,b\,\alpha')$ with
$\alpha=X_{[pd+1,\delta-1]}$ of length $\delta_0-1$ and $\alpha'=X_{[\delta+1,pd+d-1]}$ of
length $d-\delta_0-1$, the marked letter at position $\delta_0$; and
$T^{(b)}_{(k-1)d+\delta_0}=\tau(\beta\,b\,\beta')$ with
$\beta=X_{[(k-1)d+\delta_0+1,\,m-1]}$ of length $d-\delta_0-1$ and
$\beta'=X_{[m+1,\,(k-1)d+\delta_0+d-1]}$ of length $\delta_0-1$, the marked letter at
position $d-\delta_0$.

(a) If $\tau$ satisfies (2B$^{**}_d$) then $T^{(0)}T^{(0)}=T^{(1)}T^{(1)}$ in
\eqref{eq:reduced} --- the four blocks $\alpha,\alpha',\beta,\beta'$ have exactly the lengths
required by \eqref{eq:2Bstarstar} --- so \eqref{eq:reduced} holds whatever $\Pi$ is.  Since
$\delta$ was arbitrary, $\sigma$ satisfies (2B$^*_m$).

(b) Fix $\delta_0\in[1,d-1]$ and take $p=k-1$, i.e.\ $\delta=(k-1)d+\delta_0$; this is
legitimate since $\delta\le (k-1)d+d-1=m-1$.  Then the two marked windows are
$\bigl((k-1)d,\,kd\bigr)$ and $\bigl((k-1)d+\delta_0,\,(k-1)d+\delta_0+d\bigr)$, which
overlap in $\bigl((k-1)d+\delta_0,\,kd\bigr)$; comparing the displays above, $\beta=\alpha'$,
so \eqref{eq:reduced} is an instance of \eqref{eq:2Bstar} multiplied by $\Pi$.  The
remaining factors are $T_{id}$ for $0\le i\le k-2$, supported on the disjoint windows
$(id,(i+1)d)\subseteq(0,(k-1)d)$, and $T_{\delta+id}$ for $1\le i\le k-1$, supported on the
disjoint windows $(\delta+id,\delta+(i+1)d)\subseteq(kd+\delta_0,\ \delta+m)$.  All $2k-2$
of these windows are pairwise disjoint and disjoint from $\bigl((k-1)d,\,kd+\delta_0\bigr)$,
which carries $\alpha,\alpha'=\beta,\beta'$ and the marked sites.  Choose $u^{+}$ with
$\tau(u^{+})\ne0$ and set $X$ equal to $u^{+}$ on each of the $2k-2$ windows; then
$\Pi=\tau(u^{+})^{2k-2}\ne0$ because the ring is a domain.  Cancelling this one named
element from \eqref{eq:reduced} gives \eqref{eq:2Bstar} at $\delta_0$ for $\tau$, with
$\alpha,\alpha'=\beta,\beta'$ free.  As $\delta_0$ was arbitrary, $\tau$ satisfies
(2B$^*_d$).
\end{proof}

\begin{remark}[both implications of Lemma~\ref{lem:descent} are strict]\label{rem:gap}
Part (b) yields only (2B$^*_d$), not (2B$^{**}_d$), because for $p\le k-2$ the block $\beta$
sits inside one of the $\Pi$-windows, so ``$\Pi\ne0$'' and ``$\beta$ arbitrary'' cannot be
arranged at once.  That obstruction is real, not an artefact of the proof: \emph{neither
implication of Lemma~\ref{lem:descent} can be reversed.}

For $d\le3$ the conditions (2B$^{*}_d$) and (2B$^{**}_d$) have the same solutions, but for
$d\ge4$ the second is strictly stronger (Check~V5$'$: at $d=4$ they admit $6$ resp.\ $2$
proper supports, at $d=5$, $24$ resp.\ $8$).  And at $d=4$, $k=2$ there are exactly four
weights over $\F_3$, up to scaling, separating the middle condition from the strong one
(Check~V5): the singletons
\[
  \tau=\mathbb{1}_{w},\qquad w\in\{100,\ 110,\ 001,\ 011\}.
\]
Each satisfies (2B$^{*}_4$) but not (2B$^{**}_4$), and $\tau^{\otimes2}$ nevertheless
satisfies (2B$^{*}_8$) --- as does $\tau^{\otimes3}$ against (2B$^{*}_{12}$), all
$2\,096\,128$ instances.  So for $d\ge4$ and $k\ge2$,
\[
   \bigl\{\tau \text{ with (2B}^{**}_d)\bigr\}\;\subsetneq\;
   \bigl\{\tau:\ \tau^{\otimes k}\text{ has (2B}^{*}_{kd})\bigr\}\;\subsetneq\;
   \bigl\{\tau \text{ with (2B}^{*}_d)\bigr\},
\]
both inclusions strict.  Consequently Theorem~\ref{thm:counter}, which routes through
(2B$^{**}_d$), produces \emph{some} but not all commuting pairs with zeros --- the four
singletons give further ones at $e=8$ --- and Theorem~\ref{thm:main0} covers those because
it is stated with (2B$^{*}_e$) on $\sigma$, exactly what the commutator delivers.  Nothing
in Theorem~\ref{thm:main0}, Corollary~\ref{cor:main0}, Theorem~\ref{thm:gcdle2} or
Theorem~\ref{thm:counter} depends on identifying the middle condition; see
§\ref{sec:gaps}, (H1).
\end{remark}

\section{The classification and the corrected nonvanishing statement}\label{sec:main}

\begin{proof}[Proof of Theorem~\ref{thm:main0}]
By \cite[Lemma 4.1]{Q149} the commutator vanishes iff \eqref{eq:1B} and \eqref{eq:2B} vanish
for all admissible arguments.

$(\Rightarrow)$  By Proposition~\ref{prop:1B} there is a nonzero rank-$(f-e)$ weight $\rho$
over $F$ with $\sigma\otimes\rho=\bar\sigma=\rho\otimes\sigma$.  As
$\gcd(e,f-e)=\gcd(e,f)=d$, Theorem~\ref{thm:structure0} gives a nonzero rank-$d$ weight
$\tau$ over $F$ and $c,c'\in F^\times$ with $\sigma=c\,\tau^{\otimes(e/d)}$ and
$\rho=c'\,\tau^{\otimes((f-e)/d)}$; then
$\bar\sigma=\sigma\otimes\rho=cc'\,\tau^{\otimes(f/d)}$.  If $e\ge2$,
Proposition~\ref{prop:2B} gives (2B$^*_e$) for $\sigma$; if $e=1$ there is nothing to prove.

$(\Leftarrow)$  Given the stated form, set $\rho:=\bar c\,c^{-1}\tau^{\otimes((f-e)/d)}$.
Associativity of $\otimes$ gives $\sigma\otimes\rho=\bar c\,\tau^{\otimes(f/d)}=\bar\sigma
=\rho\otimes\sigma$, so \eqref{eq:1B} vanishes by Proposition~\ref{prop:1B}; and
\eqref{eq:2B} vanishes by Proposition~\ref{prop:2B} since $\sigma$ satisfies (2B$^*_e$).
\end{proof}

\begin{proof}[Proof of Corollary~\ref{cor:main0}]
$\sigma(u)=c\prod_{i=1}^{e/d}\tau\bigl(u_{[(i-1)d+1,\,id-1]}\bigr)$ and similarly for
$\bar\sigma$.  If $\tau$ vanishes nowhere then, $F$ being a field, neither does $\sigma$ or
$\bar\sigma$.  Conversely if $\tau(w)=0$ for some $w\in\{0,1\}^{d-1}$ then any $u$ whose
first $d-1$ letters are $w$ has $\sigma(u)=0$, and likewise for $\bar\sigma$.  (For $d=e$ the
first statement needs $e/d=1$; the argument is the same.)  Finally $W=\pm\sigma$ and
$\bar W=\pm\bar\sigma$ pointwise.  The two cases $\gcd\le2$ and $\gcd\ge3$ are
Theorem~\ref{thm:gcdle2} and Theorem~\ref{thm:counter} below.
\end{proof}

\subsection{$\gcd(e,f)\le2$: nonvanishing holds}

\begin{theorem}\label{thm:gcdle2}
Let $1\le e<f$ with $d=\gcd(e,f)\le2$, let $K$ be an integral domain, and let
$W\not\equiv0\not\equiv\bar W$ be shape weights of ranks $e,f$ over $K$ with
$[R_e^{W},R_f^{\bar W}]=0$.  Then no value of $W$ or of $\bar W$ is zero.
\end{theorem}

\begin{proof}
By Theorem~\ref{thm:main0} write $\sigma=c\,\tau^{\otimes(e/d)}$,
$\bar\sigma=\bar c\,\tau^{\otimes(f/d)}$ with $\tau\not\equiv0$ of rank $d$.  By
Corollary~\ref{cor:main0} it is enough to show $\tau$ vanishes nowhere.

If $d=1$ then $\tau$ is a shape weight of rank $1$, i.e.\ a single scalar $t$, and
$\tau\not\equiv0$ means $t\ne0$.  There is nothing to vanish.  (Note that this case uses
\emph{only} the one-bead sector.)

If $d=2$ then $\tau:\{0,1\}\to F$ and $d\mid e$ forces $e\ge2$, so the two-bead sector is
nonempty and Proposition~\ref{prop:2B} applies: $\sigma$ satisfies (2B$^*_e$).  By
Lemma~\ref{lem:descent}(b), $\tau$ satisfies (2B$^*_2$).  The only instance of
\eqref{eq:2Bstar} at $m=2$ is $\delta=1$ with $\pi,\pi',\varsigma$ all empty, namely
\[
   \tau(0)^2\;=\;\tau(1)^2 .
\]
If $\tau(0)=0$ then $\tau(1)^2=0$, hence $\tau(1)=0$ because $F$ is a field, hence
$\tau\equiv0$ --- a contradiction.  Symmetrically $\tau(1)\ne0$.
\end{proof}

\begin{corollary}[the coprime headline, unconditionally]\label{cor:coprime}
Let $\gcd(e,f)=1$ and $[R_e^{W},R_f^{\bar W}]=0$ with $W\not\equiv0\not\equiv\bar W$ over an
integral domain.  Then
\[
   W(u)=\alpha\,(-1)^{|u|},\qquad \bar W(y)=\beta\,(-1)^{|y|}
\]
for some $\alpha,\beta\ne0$; that is, $R_e^{W}=\alpha\,\Mult{p_e}$ and
$R_f^{\bar W}=\beta\,\Mult{p_f}$.  In particular, for $(e,f)=(1,3)$ the cross-rank
commutator forces $\bar W_{(2,1)}=\bar W_{(1,2)}$ and then the full alternation.
\end{corollary}

\begin{proof}
$d=1$, so $\tau$ is a scalar $t\ne0$ and $\tau^{\otimes k}$ is the constant function $t^{k}$;
hence $\sigma$ and $\bar\sigma$ are nonzero constants, which is the assertion.  This is
\cite[Cor.~5.3]{Q149}, now with no appeal to \cite[Lemma 8.2]{Q149} --- and, remarkably, with
no appeal to the two-bead sector either.
\end{proof}

\begin{remark}
Corollary~\ref{cor:coprime} also reproves \cite[Prop.~8.5]{Q149} (\textsf{prop:e1}, the case
$e=1$) in one line, and it does so in the regime where the two-bead sector is empty, which is
exactly the regime \textsf{prop:e1} was written to handle by hand.
\end{remark}

\subsection{$\gcd(e,f)\ge3$: nonvanishing fails}

\begin{convention}\label{conv:C}
For $d\ge3$, $1\le\delta\le d-1$ with $2\delta\ne d$, and $b\in\{0,1\}$ put
\[
   \Cset_{\delta,b}\;:=\;\bigl\{\,w\in\{0,1\}^{d-1}\ :\ w_\delta=b,\ w_{d-\delta}=1-b\,\bigr\},
\]
a nonempty set of size $2^{d-3}$ (it is nonempty precisely because $\delta\ne d-\delta$),
and let $\tau_{\delta,b}$ be its indicator function, a shape weight of rank $d$.
\end{convention}

\begin{lemma}\label{lem:Ctau}
$\tau_{\delta,b}$ satisfies \textnormal{(2B$^{**}_d$)}, hence also
\textnormal{(2B$^{*}_d$)}.
\end{lemma}

\begin{proof}
Write $\tau=\tau_{\delta,b}$, $\Cset=\Cset_{\delta,b}$.  Since $\tau$ is $\{0,1\}$-valued,
\eqref{eq:2Bstarstar} at $\delta_0$ says
\[
   \bigl[\,u^{0}\in\Cset\ \text{and}\ v^{0}\in\Cset\,\bigr]
   \quad\Longleftrightarrow\quad
   \bigl[\,u^{1}\in\Cset\ \text{and}\ v^{1}\in\Cset\,\bigr],
\]
where $u^{c}=\alpha\,c\,\alpha'$ carries $c$ in position $\delta_0$ and
$v^{c}=\beta\,c\,\beta'$ carries $c$ in position $d-\delta_0$, the other letters being
fixed.  Three cases.

If $\delta_0=\delta$: then $u^{c}\in\Cset$ requires $u^c_\delta=b$, i.e.\ $c=b$; and
$v^{c}$ carries $c$ in position $d-\delta_0=d-\delta$, so $v^{c}\in\Cset$ requires
$c=1-b$.  The two demands are incompatible, so both sides of the equivalence are false.

If $\delta_0=d-\delta$: symmetrically, $u^{c}\in\Cset$ requires $c=1-b$ (position
$d-\delta$) and $v^{c}\in\Cset$ requires $c=b$ (position $d-\delta_0=\delta$); again both
sides are false.

If $\delta_0\notin\{\delta,d-\delta\}$: the marked position of $u^{c}$ is $\delta_0\ne
\delta,d-\delta$, so membership $u^{c}\in\Cset$ is decided by the fixed letters and does not
depend on $c$; and the marked position of $v^{c}$ is $d-\delta_0$, which differs from
$\delta$ (else $\delta_0=d-\delta$) and from $d-\delta$ (else $\delta_0=\delta$), so
$v^{c}\in\Cset$ likewise does not depend on $c$.  Both sides are then equal.
\end{proof}

\begin{theorem}\label{thm:counter}
Let $1\le e<f$ with $d=\gcd(e,f)\ge3$, let $K$ be any integral domain, and let
$\tau=\tau_{1,b}$ of Convention~\ref{conv:C}.  Put
\[
   \sigma := \tau^{\otimes(e/d)},\qquad \bar\sigma := \tau^{\otimes(f/d)},\qquad
   W(u):=(-1)^{|u|}\sigma(u),\quad \bar W(y):=(-1)^{|y|}\bar\sigma(y).
\]
Then $W\not\equiv0\not\equiv\bar W$ and $[R_e^{W},R_f^{\bar W}]=0$, while $W$ vanishes at
every $u\in\{0,1\}^{e-1}$ one of whose blocks $u_{[(i-1)d+1,\,id-1]}$ lies outside
$\Cset_{1,b}$ --- that is, at $2^{e-1}-2^{(e/d)(d-3)}\cdot 2^{(e/d)-1}$ of the $2^{e-1}$
words, a positive number; likewise for $\bar W$.  In
particular \cite[Lemma 8.2]{Q149} is false for every such $(e,f)$, and with it
\cite[Theorem 5.2]{Q149}\,(necessity).  This is a sufficiency statement only: it
exhibits counterexamples, and by Remark~\ref{rem:gap} not all of them.
\end{theorem}

\begin{proof}
$\tau\not\equiv0$ since $\Cset_{1,b}\ne\emptyset$ (here $d\ge3$ gives $1\ne d-1$), so
$\sigma$ and $\bar\sigma$ are not identically zero.  By Lemma~\ref{lem:Ctau} $\tau$ satisfies
(2B$^{**}_d$), so by Lemma~\ref{lem:descent}(a) $\sigma$ satisfies (2B$^{*}_e$).  By
Theorem~\ref{thm:main0} (direction $\Leftarrow$, with $c=\bar c=1$) the commutator vanishes.
Finally $\Cset_{1,b}\ne\{0,1\}^{d-1}$, because $d\ge3$ means the constraint
$w_1=b,\ w_{d-1}=1-b$ excludes e.g.\ the word $w$ with $w_1=w_{d-1}=b$; so $\tau$ has a
zero, and Corollary~\ref{cor:main0} transports it to $W$ and $\bar W$.

That this contradicts \cite[Lemma 8.2]{Q149} is immediate.  It contradicts
\cite[Thm.~5.2]{Q149}(necessity) as well: case (c) there asserts
$W=\alpha W_e^{\gamma}$ with $\gamma$ valued in $F^\times$, and every such $W$ is nowhere
zero, while (a) and (b) are excluded by $W\not\equiv0\not\equiv\bar W$.
\end{proof}

\begin{theorem}[the smallest case, explicitly]\label{thm:counter36}
Let $K$ be any integral domain and let $W,\bar W$ be the shape weights of ranks $3$ and $6$
given in the ribbon-shape indexing of \cite[Lemma 1.3]{Q149} by
\[
  W_{(2,1)}=-1,\qquad W_{(3)}=W_{(1,1,1)}=0;
\]
\[
  \bar W_{(2,3,1)}=+1,\qquad \bar W_{(2,1,2,1)}=-1,\qquad
  \bar W_{\alpha}=0\ \text{ for the other } 30 \text{ compositions } \alpha\models6 .
\]
Then $[R_3^{W},R_6^{\bar W}]=0$, and both operators are nonzero.
\end{theorem}

\begin{proof}
This is Theorem~\ref{thm:counter} with $d=e=3$, $f=6$, $b=1$: then
$\Cset_{1,1}=\{w\in\{0,1\}^{2}: w_1=1,\ w_2=0\}=\{(1,0)\}$, so $\tau=\sigma$ is the indicator
of the occupancy word $(1,0)$, whose composition is
$\alpha(1,0)=(3-1,\,1-0)=(2,1)$ by \cite[Lemma 1.3]{Q149}; and $W(1,0)=(-1)^{1}\cdot1=-1$.
For $\bar\sigma=\tau^{\otimes2}$ the support is $\{y\in\{0,1\}^{5}: (y_1,y_2)=(y_4,y_5)=(1,0)\}
=\{(1,0,0,1,0),\,(1,0,1,1,0)\}$, with $\bar\sigma=1$ there; the corresponding compositions
are $(2,3,1)$ and $(2,1,2,1)$, and $\bar W=(-1)^{|y|}\bar\sigma$ gives $+1$ and $-1$.
Nonvanishing of the operators is \cite[Lemma 3.1]{Q149}: both ribbon shapes occur.
\end{proof}

\begin{remark}[why this one commutes, informally]
$R_3^{W}$ performs the local move $1100\mapsto0101$ on the Maya diagram: it needs
$b,b+1\in M$ and $b+2,b+3\notin M$.  The support of $\bar W$ is exactly the set of
configurations from which the \emph{same bead} can travel $b\mapsto b+3\mapsto b+6$, in one
of the two possible orders --- the two orders being separated by the occupancy of the site
$b+3$, which is the free letter $y_3$.  So $R_6^{\bar W}$ is the ``one-bead part of
$(R_3^{W})^2$'', up to the sign bookkeeping, and $\otimes$ is the operation that builds it.
Theorem~\ref{thm:structure0} says this is the \emph{only} way a commuting partner can arise.
\end{remark}

\subsection{Where the old classification was right}

\begin{proposition}\label{prop:nowherezero}
Let $\tau$ be a nowhere-zero shape weight of rank $m$ over a field $F$ satisfying
\textnormal{(2B$^*_m$)}.  Then there are $\gamma_1,\dots,\gamma_{m-1}\in F^{\times}$ with
\[
  \tau(w)=\tau(0^{m-1})\prod_{j=1}^{m-1}\gamma_j^{\,w_j},\qquad
  \gamma_\delta\,\gamma_{m-\delta}=1\quad(1\le\delta\le m-1).
\]
\end{proposition}

\begin{proof}
For $1\le j\le m-1$ let $\theta_j(w):=\tau(w\,|\,w_j=1)/\tau(w\,|\,w_j=0)\in F^{\times}$, a
function of the letters $w_i$, $i\ne j$.  Dividing \eqref{eq:2Bstar} at $\delta$ by its
left-hand side gives
\[
   \theta_\delta(\pi\,\cdot\,\pi')\cdot\theta_{m-\delta}(\pi'\,\cdot\,\varsigma)=1 ,
\]
where the first factor is a function of $(\pi,\pi')$ only and the second of
$(\pi',\varsigma)$ only.  Fix $\pi'$ and vary $\pi$: the second factor cannot change, so the
first is constant in $\pi$; varying $\varsigma$ shows the second is constant in $\varsigma$.
Hence $\theta_\delta(w)$ depends only on $w_{[\delta+1,m-1]}$, and $\theta_{m-\delta}(w)$
depends only on $w_{[1,m-\delta-1]}$.  Replacing $\delta$ by $m-\delta$ in the latter (also
in $[1,m-1]$) gives: $\theta_{\delta}(w)$ depends only on $w_{[1,\delta-1]}$.  The two index
sets $[\delta+1,m-1]$ and $[1,\delta-1]$ are disjoint, so $\theta_\delta$ is a constant
$\gamma_\delta$.  Then \eqref{eq:2Bstar} reads $\gamma_\delta\gamma_{m-\delta}=1$, and
flipping the letters of $0^{m-1}$ one at a time gives the product formula.
\end{proof}

\begin{corollary}\label{cor:agree}
In Theorem~\ref{thm:main0}, the solutions with $\tau$ nowhere zero are exactly
\cite[Thm.~5.2(c)]{Q149}: $\sigma=\alpha\,\sigma^{\gamma}_e$, $\bar\sigma=\beta\,
\sigma^{\gamma}_f$ with $\gamma:\Z/d\to F^\times$, $\gamma(0)=1$,
$\gamma(\delta)\gamma(-\delta)=1$.
\end{corollary}

\begin{proof}
By Lemma~\ref{lem:descent}(b), $\tau$ satisfies (2B$^*_d$); being nowhere zero,
Proposition~\ref{prop:nowherezero} applies at $m=d$ and gives
$\tau(w)=\tau(0^{d-1})\prod_j\gamma_j^{w_j}$ with $\gamma_\delta\gamma_{d-\delta}=1$.  Then
\[
  \sigma(u)=c\prod_{i=1}^{e/d}\tau\bigl(u_{[(i-1)d+1,\,id-1]}\bigr)
   =\sigma(0^{e-1})\prod_{k=1}^{e-1}\gamma_{k\bmod d}^{\,u_k},
\]
with the convention $\gamma_0:=1$ (the letters in positions $d,2d,\dots$ are the ones
$\otimes$ ignores), and likewise for $\bar\sigma$.  This is exactly
\cite[Conv.~5.1]{Q149}.  Conversely every such pair has nowhere-zero $\tau$.
\end{proof}

So \cite[Thm.~5.2]{Q149} is not wrong about the solutions it lists; it is incomplete.  Its
list is the nowhere-zero locus, which is everything precisely when $d\le2$.

\section{Verification}\label{sec:ver}

All code is at \texttt{proofs/code/2026-09-16-c1-Q150/}.  The numeric commutator checker
\texttt{direct.py} is written from \cite[Conv.~1.2]{Q149} and shares no algebra with
\texttt{solver.py} of \cite{Q149}.

\paragraph{V0 (calibration).}  On $(e,f)=(2,3)$, \texttt{direct.py} and
\texttt{solver.commutator\_eqs\_window} agree on all $300$ random weight pairs over
$\F_3$; the population is live, $16$ commuting and $284$ not.  A first version of
\texttt{direct.py} silently dropped the legal moves originating below the frozen tail; it
was corrected before use.  A first run compared residuals over $\Z$ against solutions over
$\F_3$ and reported four false failures (residuals $\pm3,\pm6$); the checker now takes a
modulus.

\paragraph{V0$'$ (the counterexample, two code paths).}  For
Theorem~\ref{thm:counter36}: \texttt{direct.py} finds $0$ nonzero commutator matrix
elements; \texttt{solver.commutator\_eqs\_window(3,6,Lw=14)} with the same numeric weights
sweeps $16384$ states, $4192$ matrix elements, and returns $0$ nonzero residuals over $\Z$.
Positive controls at $(3,6)$: Murnaghan--Nakayama commutes; the $\gamma$-character with
$d=3$, $\gamma(1)=5$, $\gamma(2)=1/5$ commutes.  Negative control: perturbing one value of
that $\bar W$ by $1$ makes the commutator nonzero.

\paragraph{V1 (structure theorem).}  For $(e,f)\in\{(2,3),(2,4),(3,4),(3,5),(3,6),(4,5),
(4,6)\}$, every solution of the reduced system over $\F_3$ --- $4,8,4,4,16,4,8$ solutions
respectively --- is of the form $\sigma=c\,\tau^{\otimes(e/d)}$,
$\rho=c'\tau^{\otimes((f-e)/d)}$ for a rank-$d$ $\tau$.  Zero exceptions.

\paragraph{V2 (counterexamples).}  For $(e,f)\in\{(3,6),(3,9),(6,9),(4,8),(5,10),(6,12),
(9,12),(6,15)\}$ and $\tau=\tau_{1,b}$: $0$ violations of (2B$^*_e$) and $0$ of
\eqref{eq:comm}; and for $e+f\le15$ the direct commutator is zero.  The zero counts are
recorded and factor as predicted, e.g.\ at $(6,9)$: $W$ has $30$ zeros out of $32$
($2=2^{6-1-4}$ survivors, $\Cset_{1,1}$ having size $2^{3-3}=1$ per block and two blocks),
$\bar W$ has $252$ out of $256$.
\paragraph{V3 ($\gcd\le2$).}  Exhaustively over $\F_5$ for
$(e,f)\in\{(2,3),(2,4),(2,5),(3,4),(3,5),(3,7),(4,5),(4,6),(4,7),(5,6),(4,10)\}$: no solution
has a vanishing weight value.

\paragraph{V4 (negative controls, all must fire).}
(a) lifting a rank-$3$ $\tau$ to $(e,f)=(6,10)$, where $3\nmid10$, violates \eqref{eq:comm};
(b) $\tau=\mathbb{1}_{(0,0)}$, which fails (2B$^*_3$), lifted to $(6,9)$ violates
(2B$^*_6$); (c) the direct commutator for (b) is nonzero.  All three fire.

\paragraph{V5 (descent).}  Over $\F_3$, all nonzero rank-$d$ weights $\tau$, for
$d\in\{2,3,4\}$ and $k\in\{1,2,3\}$ with $2^{kd-1}\le256$.  At $k=1$ the middle condition is
(2B$^*_d$) itself ($8/8$, $80/80$, $6560/6560$ weights at $d=2,3,4$).  At $k\ge2$,
(2B$^{**}_d$) implies (2B$^*_{kd}$) with no exception, as Lemma~\ref{lem:descent}(a)
requires; the converse holds at $d=2$ ($8/8$) and $d=3$ ($80/80$) but \textbf{fails at
$d=4$}: $6552$ of $6560$ weights agree, and the eight that do not are the four singletons
$\mathbb{1}_{100},\mathbb{1}_{110},\mathbb{1}_{001},\mathbb{1}_{011}$ and their doubles,
which satisfy (2B$^*_8$) after lifting while failing (2B$^{**}_4$).  A separate run confirms
the same four satisfy (2B$^*_{12}$) at $k=3$ (all $2\,096\,128$ instances).  This is
Remark~\ref{rem:gap}, and I record how it surfaced: the check was written expecting
agreement, and its assertion caught a claim that had already been written into the draft.

\paragraph{V5$'$ (the two two-bead conditions differ).}  Enumerating all admissible supports:
at $d=2$ both conditions admit only the full support; at $d=3$ both admit the two proper
supports $\{(0,1)\},\{(1,0)\}$; at $d=4$, (2B$^*$) admits $6$ proper supports but
(2B$^{**}$) only $2$; at $d=5$, $24$ against $8$.  This is the content of
Remark~\ref{rem:gap}.

\paragraph{Scope, stated against the quantifier.}  The exhaustive scans V1 and V3 range over
$\F_p$ for $p\in\{3,5\}$ and over the pairs listed; they are not a proof for other $(e,f)$
or other characteristics, and they are not needed for one --- §§\ref{sec:reduce}--\ref{sec:main}
are unconditional.  One earlier scan of mine \emph{was} too narrow and gave a false
negative: searching only \emph{singleton}-supported $\sigma$ reported ``no zero-bearing
solution'' at $(6,9)$, $(9,12)$, $(8,12)$, $(10,15)$ --- all with $\gcd\ge3$ --- and
suggested the criterion ``$e\mid f$''.  The composite $\tau^{\otimes2}$ at $(6,9)$ refutes
that criterion.  The reported counts had been $12,104,48,192$ candidate words passing
(2B$^*$); the quantifier those counts answered was ``how many \emph{singletons}'', not ``how
many weights''.

\section{What is and is not proved}\label{sec:gaps}

\subsection*{Proved}
\begin{enumerate}
\item Proposition~\ref{prop:1B}: the one-bead sector is exactly
  $\bar\sigma=\sigma\otimes\rho=\rho\otimes\sigma$.  Unconditional; no nonvanishing.
\item Proposition~\ref{prop:2B}: given that, the two-bead sector is exactly (2B$^*_e$) for
  $\sigma$ alone.  $\bar\sigma$ is eliminated entirely.
\item Theorem~\ref{thm:structure0}: commuting weights are powers of a common rank-$\gcd$
  weight.  Unconditional.
\item Lemma~\ref{lem:descent}: (2B$^{**}_d$) $\Rightarrow$ (2B$^{*}_m$) $\Rightarrow$
  (2B$^{*}_d$) for $\sigma=\tau^{\otimes(m/d)}$.
\item Theorem~\ref{thm:main0} and Corollary~\ref{cor:main0}: the classification and the
  reduction of nonvanishing to nonvanishing of $\tau$.
\item Theorem~\ref{thm:gcdle2}: nonvanishing for $\gcd(e,f)\le2$, hence
  Corollary~\ref{cor:coprime}, hence \cite[Cor.~5.3, Prop.~8.5]{Q149} unconditionally.
\item Theorem~\ref{thm:counter}, Theorem~\ref{thm:counter36}, Lemma~\ref{lem:Ctau}:
  explicit counterexamples for every $\gcd(e,f)\ge3$.
\item Corollary~\ref{cor:agree}: the nowhere-zero part of the new classification is exactly
  \cite[Thm.~5.2(c)]{Q149}.
\end{enumerate}

\subsection*{Not proved}
\begin{enumerate}
\item[(H1)] \textbf{The exact rank-$d$ form of (2B$^*_e$) when $e>d$.}  I prove
  (2B$^{**}_d$) $\Rightarrow$ (2B$^{*}_e$) $\Rightarrow$ (2B$^{*}_d$), and
  Remark~\ref{rem:gap} shows \emph{both} implications are strict for $d\ge4$: the condition
  in the middle is a genuinely new one, sitting strictly between, and I do not have a closed
  form for it.  The four singletons at $d=4$ say what a proof must accommodate --- ``for
  every $\beta$ there is $\eta$ with $\tau(\eta\beta)\ne0$'' is exactly what a singleton
  $\tau$ violates, yet its lift still satisfies (2B$^{*}_{kd}$) --- so the
  cofactor-cancellation route cannot be rescued by a better choice of witness.  Whether the
  middle condition depends on $k$ I do not know; it does not for $k\in\{2,3\}$ at $d\le4$.
  \emph{This does not affect any theorem above}: the classification is stated with
  (2B$^*_e$) on $\sigma$, which is exactly what the commutator gives.
\item[(H2)] \textbf{A classification of the rank-$d$ solutions of (2B$^{*}_d$) with zeros.}
  Computationally (V5$'$) the admissible supports at $d\le5$ are the full support and the
  subsets of the $\Cset_{\delta,b}$; I prove only that each $\Cset_{\delta,b}$ \emph{is}
  admissible (Lemma~\ref{lem:Ctau}), not that every proper admissible support is contained in
  one, nor which value assignments on such a support satisfy (2B$^*_d$).  This is now the
  interesting question: it is a finite, rank-$d$-only question, and it is what a complete
  description of the $\gcd\ge3$ locus requires.
\item[(H3)] \textbf{Shared weights and more than two ranks.}  Gaps (G3), (G4) of
  \cite{Q149} are untouched.  Note that (G4) looks easier now: requiring
  $[R_{e_i}^{W_i},R_{e_j}^{W_j}]=0$ for all pairs from a set makes all the $W_i$ powers of a
  common $\tau$ pairwise, and the question is whether the $\tau$'s can be taken equal ---
  again a commuting-words question.
\item[(H4)] \textbf{(G2) of \cite{Q149}}, the $\Z/2$ at $\delta=d/2$ being a gauge, is
  unaffected and still open.
\end{enumerate}

\section{Post-mortem: why $e\le2$ closed and $e\ge3$ could not}\label{sec:postmortem}

The brief asked for the step in the $e\le2$ proof that uses ribbon structure, on the grounds
that the counterexample template of $2405.15848$ shows that commutation alone cannot force
nonvanishing.  The answer is that there is no such step.  The $e\le2$ proof closes for a
reason that has nothing to do with ribbons and everything to do with arithmetic: at $e\le2$
one has $d=\gcd(e,f)\le2$, and a rank-$\le2$ weight $\tau$ has at most two values tied by
$\tau(0)^2=\tau(1)^2$, so a zero in one forces a zero in the other.  At $e=3$ with
$3\nmid f$ the lemma is still true, but again for the arithmetic reason $d\le2$; and at
$e=3$, $f=6$ it is false.  The checks C5 and C6 of \cite{Q149} covered
$(1,3),(1,4),(1,5),(2,3),(2,4),(3,4),(3,5)$: every one of these has $\gcd\le2$.  The smallest
pair with $\gcd\ge3$ is $(3,6)$, and it was never run.

Three further notes, recorded because they are the transferable part.
\begin{itemize}
\item \emph{The cylinder bookkeeping was the wrong object.}  Gap (G1) described the
  obstruction as ``the cylinder produced has length $e-\delta$, not $e-1$, so the induction
  does not terminate in one step''.  That is a true description of a route to a false
  theorem.  The reason the induction would not terminate is that there was nothing at the
  end of it.
\item \emph{Eliminating the larger unknown first.}  Every attempt so far reasoned about the
  pair $(\sigma,\bar\sigma)$, with $2^{e-1}+2^{f-1}$ unknowns.  Proposition~\ref{prop:1B}
  replaces $\bar\sigma$ by $\rho$ of rank $f-e$, and then the two-bead sector becomes a
  statement about $\sigma$ alone.  The classification is then visibly a question about a
  \emph{semigroup}, and the semigroup has a classical theorem.
\item \emph{A negative scan must be audited against its quantifier.}  See the last paragraph
  of §\ref{sec:ver}.  A scan over singleton supports reported ``no counterexample at
  $(6,9)$'' and would have produced the wrong theorem (``$e\mid f$'' instead of
  ``$\gcd(e,f)\ge3$'').  The composite $\tau^{\otimes2}$ is not a singleton.
\end{itemize}

\begin{thebibliography}{9}
\bibitem{Q149} Clio, \emph{Shape weights on ribbons: the commuting locus is the
  transpose-equivariant $\gcd(e,f)$-periodic characters} (Q149),
  \texttt{proofs/2026-09-15-c1-Q149-shape-weights.tex}, 15 September 2026.
  Numbering used above is that of the compiled PDF, checked against it on
  16 September 2026: Conv.~1.2 $=$ \textsf{conv:def}, Lemma 1.3 $=$ \textsf{lem:dict},
  Lemma 1.6 $=$ \textsf{lem:transpose}, Conv.~2.1 $=$ \textsf{conv:sigma},
  Lemma 3.1 $=$ \textsf{lem:realise}, Lemma 4.1 $=$ \textsf{lem:sectors},
  Lemma 4.2 $=$ \textsf{lem:onebead}, Lemma 4.3 $=$ \textsf{lem:twobead},
  Conv.~5.1 $=$ \textsf{conv:gamma}, Thm.~5.2 $=$ \textsf{thm:main},
  Cor.~5.3 $=$ \textsf{cor:coprime}, Prop.~7.1 $=$ \textsf{prop:suff},
  Lemma 8.1 $=$ \textsf{lem:flip}, Lemma 8.2 $=$ \textsf{lem:nonvanish},
  Lemma 8.4 $=$ \textsf{lem:char}, Prop.~8.5 $=$ \textsf{prop:e1}.
\bibitem{Q92} Clio, \emph{Ribbon commutators on Maya diagrams} (Q92), September 2026.
\end{thebibliography}

\end{document}
